\documentclass[11pt]{article}
\usepackage{fontspec}
\usepackage[french,english]{babel}
\usepackage{amsmath}
\usepackage{amssymb}
\usepackage{amsthm}
\usepackage{mathrsfs}
\usepackage{tikz-cd}
\usepackage{enumitem}

\newcommand{\printedpage}[1]{\marginpar{\footnotesize p.~#1}}
\newcommand{\msp}[1]{\marginpar{\footnotesize [#1]}}
\newcommand{\Fltr}{\operatorname{Fl}\operatorname{tr}}
\newcommand{\Ob}{\operatorname{Ob}}
\newcommand{\bDelta}{\boldsymbol{\Delta}}

\begin{document}

% ===================== page 1 =====================
\printedpage{}
% head: (none; title page, no page number printed)

\begin{center}
{\Large\textsf{THE DERIVATORS}}

\bigskip

\textbf{ALEXANDRE GROTHENDIECK}

\bigskip\bigskip

\textbf{Chapter VII}

\medskip

\textbf{Categories of paths (2)}
\end{center}

\newpage

% ===================== page 2 =====================
\printedpage{}
% head: (none; no page number printed)

This text was deciphered and transcribed in \LaTeXe\ by M. Künzer. It was edited by M. Künzer, J. Malgoire, and G. Maltsiniotis. The transcription is as faithful as possible to the manuscript. For the few obvious corrections, or rare comments by the editors, as well as for the original numbering of the manuscript pages, typewriter characters \texttt{[typewriter]} in brackets are used. A question mark in brackets means that one is not sure of the preceding word.

This edition is provisional. Remarks, comments and corrections are welcome. Send a message to:

\textbf{[e-mail address removed]}

\hfill G. Maltsiniotis\hspace*{3em}

\newpage

% ===================== page 3 =====================
\printedpage{1}
% head: A. Grothendieck : Dérivateurs, Ch. VII, Édition des Universités de Montpellier II et Paris VII
% foot: Version 22 décembre 2001

\msp{page 1}
Since the ``simplified'' formalism of paths that I wanted to work out in a hurry has turned out to be unbearably cumbersome, enough to make it virtually unusable, I am going to try to work out a flexible formalism, sticking to intuition, at the cost only of an additional conceptual effort.

\section*{1\quad Intervals and their variants} % [sic] leur variantes -> leurs variantes, plural agreement

We call a \textit{standard interval category} a category defined by an ordered set, whose set of elements is an interval of integers $[0,n]$ ($n\in\mathbf{N}$, we do not exclude the case $n=0$, i.e. of the one-point category $\{0\}$): % [sic] élements -> éléments, missing accent
\[
\Ob I=[0,n]\subseteq\mathbf{N},
\]
and satisfying the following hypothesis:

\begin{enumerate}
\item[a)] for every pair of two consecutive vertices $i-1$, $i$ of $I$ ($1\leq i\leq n$), there is an arrow joining them, thus either $i\longrightarrow i+1$, or $i\longleftarrow i+1$ \texttt{[rather} $i-1\longrightarrow i$\texttt{,} $i-1\longleftarrow i$\texttt{]} (this arrow is denoted by $\underline{u}_i$,
\[
(i-1)\xrightarrow{\underline{u}_i}i\qquad\quad\text{or}\qquad\quad(i-1)\xleftarrow{\underline{u}_i}i\ \ ).
\]
\item[b)] The category $I$ is generated by the set of its arrows $\underline{u}_i$ (called \textit{transition arrows}).
\end{enumerate}

The transition arrow is called \textit{direct} if it goes from $i-1$ to $i$, \textit{retrograde} in the opposite direction. If we assign to the interval $[i-1,i]$ the sign $+$ or $-$ according as $\underline{u}_i$ is direct or retrograde, we thus see that the
\msp{page 2}
standard intervals of given length $n$ correspond one-to-one to the systems of signs $+$ or $-$ on the $n$ transition intervals $[i-1,i]$ ($1\leq i\leq n$), so there are $2^n$ of them. The vertex $0$ is the \textit{origin} or the \textit{source} of the interval, the vertex $n$ the \textit{endpoint} or the \textit{target} of the interval.

A standard interval is \textit{oriented,} in that we have the natural \textit{total} order relation on the set of its vertices, namely that of the integers in $[0,n]$. When we speak of the order relation on $I$, that a vertex $i$ \textit{``precedes''} or \textit{``follows''} another one $j$ ($i\leq j$ or $i\geq j$ respectively), it is this order relation that we will mean. The one associated with the category structure will never be regarded as an order relation. There is also a natural total order relation on the set of transition arrows $\underline{u}_i$ ($1\leq i\leq n$), given by the order relation of the integers $i$, $1\leq i\leq n$, that index them. We will say that a transition arrow $i$ \texttt{[rather} $\underline{u}_i$\texttt{]} \textit{precedes} or \textit{follows} another \texttt{[}$\underline{u}_j$\texttt{]}, or that it is the \textit{immediate successor} or \textit{immediate predecessor} of another (if $i=j+1$ or $i=j-1$). % [sic] qu'on à -> qu'on a, verb; [sic] un sommet i en ‘précède’ -> un sommet i ‘précède’, superfluous "en"

% Reader layout: continuous text; edition coordinate retained below.
% ===================== page 4 =====================
\printedpage{2}
% head: A. Grothendieck : Dérivateurs, Ch. VII, Édition des Universités de Montpellier II et Paris VII
% foot: Version 22 décembre 2001

If we have a set of integers
\[
E\subseteq\mathbf{Z},
\]
we will consider in $E$ the equivalence relation generated by the relation between two elements $(i,j)$ of $E$: $j$ is the immediate successor of $i$ \textit{in} $\mathbf{Z}$ (i.e. $j=i+1$). Two distinct elements $i$, $j$ of $E$ are congruent for this relation if and only if the interval of integers that they determine is $\subseteq E$, thus % [sic] distinct -> distincts, plural agreement
\[
\text{if } i<j,\ i \text{ and } j \text{ are \textit{``connected''} in } E \text{ if and only if } [i,j]\subseteq E.
\]

\msp{page 3}
When we speak of the ``connected components'' of $E$, it will always be a matter of the classes under this equivalence relation.

We will identify the set of transition arrows $\Fltr(I)$ of $I$ with the interval of integers $[1,n]$, hence with a subset of $\mathbf{Z}$. If we have a subset $E\subseteq\Fltr(I)$, we thus have in it the notion of transition arrows ``connected'' in $E$, and the notion of connectedness for $E$. If $i<j$, $\underline{u}_i$ and $\underline{u}_j$ are connected if and only if all the intermediate transition arrows between $\underline{u}_i$ and $\underline{u}_j$ are in $E$. And $E$ is connected if and only if it is either empty, or formed of a sequence of transition arrows each of which is the immediate successor of the one preceding it, % [sic] soit formée -> soit formé, agreement with "il"
\[
E=\{\underline{u}_i,\underline{u}_{i+1}\ldots,\underline{u}_{i+l-1}\}, % [sic] \underline{u}_{i+1}\ldots -> \underline{u}_{i+1},\ldots, missing comma
\]
where $l$ is the cardinal of $E$. Such an $E$ is called a \textit{chain} of transition arrows, and $l$ is its length. A chain is called \textit{composable} if the composite $\underline{u}_{i+l-1}\circ\cdots\circ\underline{u}_{i+1}\circ\underline{u}_i$ \texttt{[or} $\underline{u}_i\circ\underline{u}_{i+1}\circ\ldots\circ\underline{u}_{i+l-1}$\texttt{]} makes sense, i.e. if they are all in the same direction, either direct, or retrograde.

We denote by $\Fltr^+(I)$, $\Fltr^-(I)$ the set of direct, or retrograde, transition arrows, respectively. Their connected components are nothing other than the maximal composable chains. Thus, the maximal composable chains form a partition of the set of transition arrows, into mutually disjoint intervals of integers. The chains themselves form a totally ordered set
\[
C_1,\ C_2,\ \ldots\ C_r\qquad\quad(1\leq r\leq n).
\]
The integer $r$ is called the \textit{``useful length''} of the interval
\msp{page 4}
$I$, it is thus the cardinal of the set of its maximal composable chains.

Two maximal composable chains \texttt{[consecutive]} are always of opposite sign, if one is direct the other is retrograde, and conversely. The first maximal composable chain $C_1$ (let us say simply \textit{``structural chain''}) has the same sign as $\underline{u}_1$; \texttt{[we]} say that the \textit{initial direction} of $I$ is direct, resp. retrograde, if this is so for $\underline{u}_1$, or what amounts to the same thing, for $C_1$. Likewise we define the \textit{final direction} of $I$ as that of $\underline{u}_n$, or again that of $C_r$. Initial and final directions are equal if the useful length $r$ is
\printedpage{3}
% head: A. Grothendieck : Dérivateurs, Ch. VII, Édition des Universités de Montpellier II et Paris VII
% foot: Version 22 décembre 2001
odd, unequal if it is even. (\textbf{NB} It is only for $n=0$, i.e. when $I=\{0\}$ is the one-point category, that initial direction and final direction are not defined.)

The composable chains correspond to the full subcategories of $I$ that are isomorphic to a category $\bDelta_i$ (where $i$ is the length of the chain). Let us note that $\bDelta_i$ has no automorphism other than the identity, so the maximal chains of given length $i\geq 1$ correspond exactly to the functors
\[
\bDelta_i\lhook\joinrel\longrightarrow\ I
\]
that are injective on the vertices. Such a functor is fully faithful, and a monomorphism. This functor defines a strictly increasing or strictly decreasing map on
\msp{page 5}
the vertices, according as the chain is direct or retrograde.

Let us note that every standard simplex $\bDelta_i$, $i\geq 0$, can be viewed as an interval, of length $i$, of useful length $1$ (if $i\geq 1$, $0$ if $i=0$). Its origin is $0$, its endpoint is $i$. But let us call an arbitrary category an \textit{interval category} if it is isomorphic to a standard interval. Let us note here

\textbf{Proposition 1.} \textit{Let $I$ be a standard interval of length $n\geq 1$.}

\begin{enumerate}
\item[\textit{a)}] \textit{Let $i\in I$. For $i\in\{0,n\}$, i.e. for $i$ to be one of the two} ``ends'' \textit{(or} ``endpoints''\textit{) (source or target) of $I$, it is necessary and sufficient that one of the two categories $I/i$ or $i\backslash I$ be the one-point category and that the other be isomorphic to a} $\bDelta_j$ (\footnote{where $j$ is the length of the unique structural chain that ``touches'' $i$ (i.e. such that $i$ is in the full subcategory defined by the said chain).}). \textit{(Intrinsic characterization of the ends of an interval.) For $i$ to be the source, it is necessary and sufficient that $i\backslash I$} \texttt{[incomplete sentence]}
\item[\textit{b)}] \textit{Every functor $I\longrightarrow I$ that fixes the two endpoints of $I$ is the identity.}
\end{enumerate}

\textbf{Corollary.} \textit{The group of automorphisms of $I$ has at most two elements. If it has two, the nontrivial automorphism interchanges the two endpoints. (We then say that the interval is} symmetric.\textit{)}

\msp{page 6}
NB If $i$ is not an end, and if nevertheless one of the two categories $i\backslash I$, $I/i$ is one-point, say $I/i$, then in the immediate neighborhood of $i$, $I$ has the form
\[
\begin{tikzcd}[row sep=0pt]
\bullet & \bullet \arrow[l, "\underline{u}_i"'] \arrow[r, "\underline{u}_{i+1}"] & \bullet \\
i-1 & i & i+1
\end{tikzcd}
\]
($\underline{u}_i$ retrograde, $\underline{u}_{i+1}$ direct - if $i\backslash I$ were the one-point one, then $\underline{u}_i$ would be direct, $\underline{u}_{i+1}$ retrograde - in any case these two arrows have opposite directions, i.e. not composable). We thus find that $i\backslash I$ is a category of the form

\end{document}
